% TODO
\section{Aplicación de la IA y critica cruzada}
\subsection*{Perspectiva de Kant:}
Para kant un agente autónomo no considera que tenga máxima, esto porque aún no ha sido demostrado que este tenga intenciones, por lo tanto no existiría la máxima, en consecuencia no se podría definir si el agente comete actos morales o inmorales usando el Imperativo Categórico. Entonces la Inteligencia Artificial para Kant no es más que una herramienta en la que estadísticos y matemáticos adjuntan leyes morales en él con intenciones predispuestas en el algoritmo o modelo. 
Para casos de índole ético que interviene la Inteligencia Artificial se debe de juzgar las intenciones y para crear un sistema o arquitectura que usa IA se debe de filtrar las posibles máximas que hacen posible que la IA logre su objetivo, para que con ellos se use los principios del Imperativo Categórico para determinar que la IA, en la toma de decisiones, haga un acto moral. 
\subsection*{}¿Qué pensaría Kant del pensamiento de Mill acerca de un algoritmo de redes sociales que prioriza contenido adictivo porque genera un alto nivel de satisfacción y entretenimiento (placer) en la mayoría de los usuarios?
\subsection*{Al verificar el acto con la intención de mejorar ganancias como fin y usando a la personas como un medio, sería por lo tanto un acto inmoral, debido a que no cumpliría con el segundo principio del Imperativo Categórico que es la Humanidad, las personas no deben ser usadas como un medio sino como un fin solamente. Es por eso que negaría el pensamiento de Mill quien establece el beneficio común.}
\subsection*{}¿Qué pensaría Kant del pensamiento de Bentham acerca de un vehículo autónomo programado para atropellar a un peatón si eso evita un choque donde morirían cuatro pasajeros?
\subsection*{Kant repudiaría la idea de tratar de utilizar a un conjunto de personas como un medio para decidir un acto inmoral, es decir estaría en contra de realizar el algoritmo que determine el acto que cometerá. Sin embargo, sin este algoritmo, Kant estaría de acuerdo en que suceda un evento trágico determinista, la cual análogamente sería otro acto de lo fenoménico.  }
\subsection*{}¿Qué pensaría Kant del pensamiento de Mill acerca de una IA de vigilancia que elimina la privacidad en una ciudad, argumentando que la seguridad resultante produce un bienestar superior para la población?
\subsection*{}
\subsection*{}¿Qué pensaría Kant del pensamiento de Mill acerca de un software de selección de personal que excluye a candidatos brillantes de grupos minoritarios porque su inclusión podría causar fricciones que reduzcan la productividad general de la empresa?
\subsection*{}
